\documentclass[12pt]{article}
\usepackage[utf8]{inputenc}
\usepackage{listings}
\usepackage{xcolor}
\usepackage{geometry}
\usepackage{parskip}

\geometry{margin=2.5cm}

\definecolor{codebg}{RGB}{245,245,245}
\definecolor{keyword}{RGB}{0,112,0}
\definecolor{comment}{RGB}{128,128,128}
\definecolor{string}{RGB}{180,0,0}
\definecolor{outputbg}{RGB}{240,248,255}

\lstdefinestyle{pythonstyle}{
    backgroundcolor=\color{codebg},
    basicstyle=\ttfamily\small,
    keywordstyle=\color{keyword}\bfseries,
    stringstyle=\color{string},
    commentstyle=\color{comment}\itshape,
    language=Python,
    frame=single,
    rulecolor=\color{gray!40},
    breaklines=true,
    showstringspaces=false,
    tabsize=4,
    captionpos=b,
}

\lstdefinestyle{outputstyle}{
    backgroundcolor=\color{outputbg},
    basicstyle=\ttfamily\small,
    frame=single,
    rulecolor=\color{blue!20},
    breaklines=true,
}

\title{\textbf{Wind Turbine Blade Fault Classification\\Using a Simple CNN}\\[0.5em]
\large Step 10 --- Evaluate the Model}
\author{Amreen Batool}
\date{}

\begin{document}
\maketitle

\section*{Step 10 --- Evaluate the Model}

\begin{lstlisting}[style=pythonstyle]
test_loss, test_accuracy = model.evaluate(test_data)

print("Test Loss:",     test_loss)
print("Test Accuracy:", test_accuracy)
\end{lstlisting}

\subsection*{Explanation}

This step checks the model's performance on the \textbf{unseen test images} — data the model has never seen during training or validation.

\begin{itemize}
    \item \texttt{test\_loss} $\rightarrow$ the model's error on the test set. Lower is better.
    \item \texttt{test\_accuracy} $\rightarrow$ the fraction of test images classified correctly.
\end{itemize}

\subsection*{Use in Report}

After running this step, note your values and write in your report:

\begin{quote}
\textit{The model achieved a test accuracy of approximately \textbf{XX\%} and a test loss of \textbf{YY}.}
\end{quote}

Replace \textbf{XX} and \textbf{YY} with your actual printed values.

\end{document}
